\section{Introduction}

Lie algebras were originally introduced by Sophus Lie in the 1870s to study the concept of continuous transformation groups. Lie's motivation was to develop a theory for differential equations analogous to Galois theory for algebraic equations. He discovered that the infinitesimal transformations of a continuous group form a linear space equipped with a bracket operation, now known as a Lie algebra. This linearization allows complex global geometric problems to be translated into algebraic problems, making them more tractable.

Let $ F $ be a field (mostly $ \operatorname{char}F =0$, even $ F=\mathbb{C} $).

\begin{definition}
    A linear algebra $ A $ is a vector space over $ F $, with a binary bilinear operation:
    \begin{itemize}
        \item binary operation: $ A \times A \to A $, $ (a,b) \mapsto ab $.
        \item bilinear: this operation is linear in each variable. In other word, $ \forall \alpha \in F $ and $ a,a',a'',b,b',b'' \in A $, we have:
              \begin{align*}
                  (\alpha a)b=a(\alpha b) = \alpha (ab) \\
                  (a'+a'')b = a'b+a''b                  \\
                  a(b'+b'')=ab'+ab''
              \end{align*}
    \end{itemize}
    The algebra $ A $ is associative if $ \forall a,b,c \in A $, $ (ab)c=a(bc) $.
\end{definition}

\begin{example}
    $ M_n(F) $, the algebra of $ n \times n $ matrices over $ F $.
\end{example}

\begin{example}
    Let $ V $ be a vector space, $ \operatorname{End}_F (V) =\{ \text{linear transformations } \varphi: V\to V \}$ is an associative algebra with the composition operation$ (\phi \psi) (v)=\phi(\psi(v))$.

    If $ \dim V < \infty $, when we fix a basis, then we fet a matrix representation of $ \varphi $ in our fix basis.
\end{example}

The most important thing everyone should know is Wedderburn's theorem.

\begin{theorem}[Wedderburn-Artin]
    Let $A$ be a semisimple Artinian ring (with 1). Then $A$ is a direct sum of matrix rings over division rings.

    Specifically, if ${}_A A \simeq S_1^{n_1} \oplus \dots \oplus S_r^{n_r}$ where $S_1, \dots, S_r$ are non-isomorphic simple modules occurring with multiplicities $n_1, \dots, n_r$, then
    \[
        A \simeq M_{n_1}(D_1) \oplus \dots \oplus M_{n_r}(D_r) \quad
    \]
    where $D_i = \operatorname{End}_A(S_i)^{op}$. This is a ring isomorphism, and each $M_{n_i}(D_i)$ is a 2-sided ideal.
\end{theorem}

\begin{proof}
    First, since $A$ is a semisimple ring, its left regular module ${}_A A$ can be decomposed into a direct sum of simple left $A$-modules:
    \[
        {}_A A \simeq S_1^{n_1} \oplus \dots \oplus S_r^{n_r} \quad
    \]
    where $S_1, \dots, S_r$ are pairwise non-isomorphic simple modules.

    By Schur's lemma, the endomorphism ring of a simple module is a division ring, so $D_i = \operatorname{End}_A(S_i)^{op}$ is a division ring.

    Next, we compute the endomorphism ring of the left regular module $\operatorname{End}_A({}_A A)$. On one hand, for any ring $A$ with 1, we have the ring isomorphism $\operatorname{End}_A({}_A A) \simeq A^{op}$. On the other hand, using our direct sum decomposition:
    \[
        \operatorname{End}_A({}_A A) \simeq \operatorname{End}_A(S_1^{n_1}) \oplus \dots \oplus \operatorname{End}_A(S_r^{n_r}) \quad
    \]
    This decomposition holds because $\operatorname{Hom}_A(S_i^{n_i}, S_j^{n_j}) = 0$ whenever $i \neq j$.

    For each component in the direct sum, we can rewrite the endomorphism ring using matrices over $D_i$:
    \[
        \operatorname{End}_A(S_i^{n_i}) \simeq M_{n_i}(D_i)^{op} \simeq M_{n_i}(D_i^{op}) \quad
    \]
    Substituting this back into our equation for $A^{op}$, we obtain:
    \[
        A^{op} \simeq M_{n_1}(D_1^{op}) \oplus \dots \oplus M_{n_r}(D_r^{op})
    \]

    To recover the ring $A$, we take the opposite ring on both sides. Recall the property that $M_n(R)^{op} \simeq M_n(R^{op})$. Applying this to each block yields:
    \[
        A \simeq (A^{op})^{op} \simeq \left( \bigoplus_{i=1}^r M_{n_i}(D_i^{op}) \right)^{op} \simeq \bigoplus_{i=1}^r M_{n_i}(D_i)
    \]
    Thus, the isomorphism is established.
\end{proof}

\begin{remark}
    If the base field $F$ is algebraically closed, every finite-dimensional
    division algebra over $F$ is equal to $F$. Hence the theorem simplifies to
    \[
        A \cong \prod_{i=1}^{r} M_{n_i}(F).
    \]
\end{remark}

\begin{example}
    Let $G$ be a finite group and let $\mathbb{C}[G]$ be its group algebra.
    By Maschke's theorem $\mathbb{C}[G]$ is semisimple. Therefore
    \[
        \mathbb{C}[G] \cong \prod_{i=1}^{r} M_{n_i}(\mathbb{C}),
    \]
    where $n_i$ are the dimensions of the irreducible complex representations
    of $G$.
\end{example}

\begin{definition}
    Suppose $ F $ is algebraically closed. An algebra $ A $ is simple if
    \begin{enumerate}
        \item $ \nexists \  I \triangleleft A $ and $ I \neq 0 $.
        \item $ A\cdot  A \neq (0)$.
    \end{enumerate}
\end{definition}

\begin{corollary}
    Any simple finite dimensional associative algebra is isomorphic to $ M_n(F) $ for some $ n $.
\end{corollary}

\begin{definition}[Jacobson Radical]
    Let $A$ be an associative algebra with identity over a field $F$.
    The \emph{Jacobson radical} of $A$, denoted by $\operatorname{Rad}(A)$
    or $J(A)$, is defined as the intersection of all maximal left ideals of $A$:
    \[
        \operatorname{Rad}(A)
        =
        \bigcap_{M \text{ maximal left ideal of } A} M.
    \]
\end{definition}


\begin{proposition}
    Let $A$ be an associative algebra with identity. Then
    \[
        \operatorname{Rad}(A)
        =
        \bigcap_{S \text{ simple left } A\text{-module}}
        \operatorname{Ann}_A(S),
    \]
    where
    \[
        \operatorname{Ann}_A(S)
        =
        \{a\in A \mid aS=0\}
    \]
    is the annihilator of $S$ in $A$.
\end{proposition}


\begin{proposition}[Equivalent Characterization]
    Let $A$ be an associative algebra with identity.
    Then an element $a\in A$ lies in $\operatorname{Rad}(A)$
    if and only if for every $x\in A$ the element
    \[
        1-ax
    \]
    is invertible in $A$.
\end{proposition}


\begin{proposition}
    The Jacobson radical $\operatorname{Rad}(A)$ is a two-sided ideal of $A$.
\end{proposition}


\begin{definition}[Semisimple Algebra]
    An associative algebra $A$ is called \emph{semisimple} if
    \[
        \operatorname{Rad}(A)=0.
    \]
\end{definition}

